
\documentclass[12pt]{article}
\usepackage[margin=1in]{geometry}
\usepackage{setspace}
\title{Politics as Smoothing\\A DISTANCE Interpretation of Politics}
\author{Kangbin Yim\thanks{Email: yim@sch.ac.kr (Dept. of Information Security Engineering, Soonchunhyang University)}}
\date{Apr. 2026}
\begin{document}
\maketitle

We often describe politics through familiar labels: progressive and conservative, left and right.
Yet from the perspective of DISTANCE, these categories may not represent the deepest layer of political reality.
Progressivism is not the essence of politics.
Conservatism is not the essence of politics.
They may simply be manifestations of a larger smoothing process.

Throughout this work, we have explored the universe as a collection of energy islands.
Differences in potential create energy gaps.
Energy gaps generate momentum.
And momentum seeks resolution.
Politics may be one of those processes.

Human history is often described as a story of progress:
kings were replaced by democratic institutions, slavery was abolished,
women gained political rights, and citizens obtained representation.
From a DISTANCE perspective, these events share a deeper structural similarity.
Each represents the reduction of a previously existing political energy gap.

A king and a citizen are different political potentials.
A slave owner and a slave are separate political islands connected by an asymmetry of power.
Men and women in societies with unequal political participation similarly occupied different political potentials.

Whenever such differences become sufficiently large, momentum accumulates.
And accumulated momentum eventually seeks release.
Revolutions, reforms, and democratic movements can therefore be interpreted as manifestations of political smoothing.

Democracy is not merely a governmental system.
It is a process through which previously separated political islands become integrated into a broader political community.
This perspective also offers a different way of understanding conservatism and progressivism.

Conservatism often emphasizes competition, efficiency, specialization, and concentration.
Progressivism often emphasizes fairness, inclusion, redistribution, and equality.
These are not simply competing moral systems.
They represent different relationships to momentum.
Conservative structures often allow differences to emerge and accumulate.
Progressive structures often attempt to reduce those differences once they become large.

One creates local momentum.
The other dissipates it.
Without differentiation, no momentum emerges.
Without smoothing, momentum accumulates indefinitely.
Politics becomes the visible expression of this oscillation.

The fundamental question is not whether progressives or conservatives are correct.
The deeper question is:
What political islands currently exist?
What momentum exists between them?
And how is that momentum being resolved?
Politics may ultimately be understood as another name for the smoothing process occurring within human civilization.

\end{document}
